%==================================Literature Survey.tex=============================
\chapter{LITERATURE SURVEY}
\label{ch:literature_survey}

Two distinct but interrelated bodies of research contribute to the evolving landscape 
of Security Information and Event Management (SIEM) and real-time threat detection.

\section{Adaptive Thresholds and Rule-Based Detection in SIEM}
Extensive research has been conducted on intrusion detection systems (IDS) and SIEM platforms 
that employ threshold-based and rule-driven methodologies for threat identification. 
Traditional approaches rely on static thresholds and manually tuned rule sets, which, 
while effective for known attack patterns, struggle to adapt to evolving threat landscapes. 
Recent work has explored adaptive threshold adjustment using statistical methods such as 
z-score normalization and rolling percentiles, enabling dynamic sensitivity adjustment based 
on baseline traffic behavior. The correlation of security events into meaningful offenses, 
first popularized by QRadar's offense engine, has been studied extensively for its ability 
to reduce alert fatigue by grouping related incidents into coherent attack narratives. 
Research into correlation engine design has explored rule sequencing, temporal windowing, 
and the tradeoff between detection precision and recall in high-volume event streams. 
The SecuriSphere project implements a correlation engine with three matcher strategies 
(event_id-based, script-based, and interval-based) that build upon these established 
design patterns while providing configurable clustering windows for offense grouping.

\section{MITRE ATT\&CK Mapping and Frameworks}
The MITRE ATT\&CK framework has become a foundational taxonomy for mapping detected 
security events to known adversary techniques and tactics. Research into automatic MITRE 
mapping of IDS alerts and SIEM events has demonstrated significant improvement in 
investigative efficiency when analysts can filter and search by technique coverage. 
Heatmap visualizations of technique coverage enable rapid identification of gaps in 
detection coverage and prioritize rule development. Studies have explored the use of 
natural language processing to extract technique identifiers from unstructured event 
descriptions, reducing manual labeling overhead. The SecuriSphere MITRE ATT\&CK panel 
provides interactive technique heatmaps and drill-down capabilities, allowing analysts 
to visualize detection coverage across the ATT\&CK matrix and identify under-monitored 
tactics in real time.

\section{UEBA and Behavior Analytics}
User and Entity Behavior Analytics (UEBA) has emerged as a critical component for 
detecting insider threats and compromised accounts through statistical anomaly detection. 
Baseline models constructed from rolling daily or hourly activity metrics are compared 
against real-time event streams, with anomalies flagged via z-score thresholding. 
Research has established that baselines requiring a minimum number of observed days 
before anomaly detection can achieve better precision-recall tradeoffs. The integration 
of UEBA into SIEM platforms bridges the gap between point-in-time alert detection and 
context-rich investigation, enabling analysts to correlate anomalous user behavior with 
broader attack timelines. The SecuriSphere UEBA module constructs rolling baselines 
for monitored entities and flags deviations using adaptive thresholding, supporting 
the detection of compromised credentials and unusual data access patterns.

\section{Correlation Engine Design and Offense grouping}
Research into correlation engine design has explored rule sequencing, temporal windowing, 
and the tradeoff between detection precision and recall in high-volume event streams. 
The offense grouping paradigm, popularized by next-generation SIEM platforms, transforms 
raw security events into actionable incidents through temporal and logical clustering. 
Key design considerations include window duration selection, de-duplication strategies, 
and the propagation of contextual metadata (such as MITRE techniques and risk scores) 
through the offense lifecycle. The SecuriSphere correlation engine implements 30-minute 
temporal clustering with script-based and event_id-based matching as alternative strategies, 
providing configurable options for different deployment environments and threat profiles.

\section{Summary and Gap Analysis}
Collectively, these studies highlight the maturity of analytical techniques—statistical 
thresholding, rule-based correlation, MITRE mapping, and UEBA—in security operations. 
However, a noticeable gap exists in accessible, open-source SIEM deployments that 
integrate rule-based detection, multi-strategy correlation, UEBA, and MITRE ATT\&CK 
visualization within a single, unified stack suitable for educational or resource-constrained 
environments. While enterprise-grade platforms (e.g., Splunk, IBM QRadar, Elastic SIEM) 
offer extensive feature sets, they require significant licensing and infrastructure investment. 
This project addresses that gap by providing a complete, Dockerized SIEM platform with 
extensible detection (13 OpenSigma rule types), three-correlation matcher algorithms, 
offense grouping with 30-minute temporal clustering, UEBA anomaly detection, and MITRE 
ATT\&CK visualization, all within a unified Next.js frontend and FastAPI backend architecture 
designed for horizontal scalability and educational deployment.

\noindent\textbf{Key Contributions of This Project:}
\begin{itemize}
    \item Full-stack SIEM architecture spanning Linux agent, FastAPI backend, and Next.js frontend.
    \item 13 OpenSigma-based detection rule types with pattern-matching extensibility.
    \item Three-strategy correlation engine (event_id, script, interval-based) with configurable clustering.
    \item 30-minute temporal offense grouping with MITRE ATT\&CK enrichment and risk scoring.
    \item UEBA module with rolling baseline construction and z-score anomaly detection.
    \item 28-page interactive dashboard covering threat intelligence, offenses, risk trends, and timelines.
    \mbox{}\newline
\end{itemize}
\end{parameter>